% 本文件由 LaTeX 排版 Agent 根据有效 translation.json 自动生成。
% author=Eusebius of Caesarea (c. 265-c. 340); work_id=2503; title=君士坦丁演说

\chapter{君士坦丁演说}

% structure: p0001 source=h1 target=omitted reason=page-title-already-rendered-by-parent

% structure: p0002 source=h2 target=section reason=relative-heading-rank; base=h2

\section{第一章 \Emph{论复活节庆期引言：论天主圣言赐予人类诸多恩惠，如何被受惠者背叛。}}

那比白昼和太阳更耀眼的光芒，复活的初果，早已分解之身体的更新，神圣应许的标记，通往永生的道路——简言之，受难之日——已经到来，至爱的师长们，以及你们，在此聚集的朋友们，有福的群众，你们敬拜那一切敬礼的创始者，并按照祂圣言的训诫，以心灵和声音不住地赞美祂。然而，你，大自然，万有之母，你曾为人类成就过何种类似这样的祝福？不，还有什么在任何意义上可算是你的作品，既然那形成宇宙者本身就是你存在的作者？因为正是祂以美装饰了你；而自然的美就是遵循自然律的生命。但与自然完全相反的原则却大大盛行，因为人们一致拒绝向万有之主献上祂应得的敬拜，相信宇宙的秩序并非依赖祂的天主眷顾，而是依赖盲目的偶然之不确定性；尽管祂所感发的先知们最清楚地宣告了真理，他们的话本应赢得信任，却遭到那恨恶真理之光、喜爱黑暗之迷途邪恶的全面抵抗。这种错误并非没有暴力和残忍相伴，特别是君王的意志助长了群众的盲目冲动，或者更确切地说，它自己在鲁莽愚蠢的道路上领路。这些原则，经过许多世代的实践，在早期时代成为可怕邪恶的源头：但救主临在的光辉一出现，正义就取代了错误，风暴的混乱之后来临了平静，先知的预言全都应验了。因为祂以自己荣耀的明辨和纯洁的品格光照了世界，并升到父家的居所之后，在地上建立了祂的教会，作为德行的圣殿，一个不朽的、永不朽坏的圣殿，在其中应虔诚地举行对至高圣父及对祂自己的敬拜。然而，各国疯狂的恶意在这里图谋什么？他们竭力拒绝基督的恩宠，并摧毁那为所有人的救恩而设立的教会，尽管他们由此确保了自身迷信的覆灭。于是，不圣洁的叛乱，战争与纷争再次盛行，伴随着刚愎、奢侈的放纵，以及对财富的渴求——这种渴求时常用虚假的希望安抚其受害者，时而用无端的恐惧打击他们；这是一种违反自然的渴求，正是邪恶本身的特性。让她（邪恶）卧于尘土，承认德行的胜利力量；让她在忏悔的苦痛中撕裂自己，正如她所能做的。但现在让我们继续谈论属于神圣教义的主题。\par

% structure: p0004 source=h2 target=section reason=relative-heading-rank; base=h2

\section{第二章 \Emph{向教会及听众呼吁，请求宽恕并纠正其演讲中的错误。}}

请听，你这船舶的舵手，拥有童贞纯洁者，以及你，教会，温柔与未经历世故年代的抚育者，真理与温良的守护者，通过你那永不停息的泉源，救恩的溪流涌出！也请宽容，我的听众们，你们真诚敬拜天主，因此是祂眷顾的对象：请留意，不是所说的言辞，而是所说的真理；不是说话的人，而是那使他的讲论神圣化的虔诚热忱！因为当说话者的真正意图仍不明时，言辞有何用呢？的确，我可能尝试谈论伟大的事；那充满我灵魂的天主之爱，一种超越自然矜持的爱，是我大胆尝试的托辞。因此，我向你们呼吁，你们这些最深谙天主奥秘的人，请以你们的建议帮助我，以你们的思绪追随我，并纠正我言辞中任何有错误之处，不求完美知识的展示，而是仁慈地接纳我努力的真挚。愿圣父和圣子的圣神赐予祂大能的助佑，当我说出祂所启示于言语或思想的话语时。因为如果有人，无论是在演讲艺术或其他任何技艺中，期望在没有天主帮助的情况下完成一部作品，那么作者和他的努力都将被发现同样不完美；而一旦蒙受上天灵感的祝福，他就没有理由恐惧，没有空间沮丧。因此，请求你们宽恕这篇序言的冗长，让我们以最大的范围尝试这个主题。\par

% structure: p0006 source=h2 target=section reason=relative-heading-rank; base=h2

\section{第三章 \Emph{论天主是圣言的圣父，万有的创造者；以及若万物的原因多样，物质对象便不能持续存在。}}

天主，祂永远超越万有，是万物所渴望的美善，祂无起源，故无开端，祂自身就是一切受造物的根源。但由祂而出者又回归于祂；这种分离与结合并非空间性的，而是理智性的。因为这种生出并不伴随圣父的实体的减损（如同由种子生出那样）；而是藉着预知的决定行动，天主显明了一位救主，掌管这个可感的世界及其中的一切受造物。因此，万有在此世界范围内的一切存在和生命的源头由此而来；灵魂和一切感官由此而来；那些感官知觉得以完成的器官也由此而来。那么，这个论证的目的是什么？为了证明万有存在着一位管理者，并且一切事物，无论天上或地上的，无论是自然体还是有机体，都服从于祂单一的统治。因为如果这些无数的事物的统治权不是掌握在一人手中，而是掌握在多人手中，就必然会有元素的划分和分配，古老的寓言就会成真；此外，嫉妒和野心，争夺更高的权力，会破坏整体的和谐一致，而每位主宰者会以不同于他人的方式管理其所控制的部分。然而，这个宇宙秩序始终如一的事实，证明它受到一位更高力量的照顾，并且其起源不能归因于偶然。否则，宇宙本性的作者如何能被认识？祈祷和恳求应首先向谁、最后向谁呈献？我该选择谁作为我崇拜的对象，而不至于对其他神明犯下不敬之罪？再者，如果我偶然想要获得某种暂时的祝福，我是否会在向那成全我请求的力量表达感谢时，向那反对它的力量传达责备？或者，当我想要知道灾难的原因并寻求解脱时，我该向谁祈祷？或者，假设神谕和预言给出了答案，但此事不在其权柄范围内，属于另一位神祇的领域。那么，怜悯在哪里？天主对人的眷顾在哪里？除非，确实有一些更仁慈的力量，对另一位没有这种感受的神祇采取敌对态度，愿意给予我保护。由此，愤怒、不和、相互指责，最终导致普遍的混乱，而每位神祇都离开了自己适当的行动范围，因野心勃勃地贪爱权力而对分配给自己的份额不满。那么，这些事情的结果会是什么？这些天上力量之间的不和肯定会对地上的利益造成破坏：时间和季节的有序更替将消失；大地的连续产出将不再被享用；白昼本身和随之而来的夜晚的安息也将停止。但这个话题已足够：让我们再次回到那种无可反驳的推理方式。\par

% structure: p0008 source=h2 target=section reason=relative-heading-rank; base=h2

\section{第四章 \Emph{论偶像崇拜的错误}}

凡有开端者，必有终结。从时间上说，开端被称为受生；凡由受生而来的，必腐朽，其美丽会随时间流逝而减损。那么，源于可腐朽之受生者，如何能不死？此外，无知大众相信这样一种假设：婚姻和生儿育女在众神之间是常见的。那么，即使这样的后代是不朽的，且不断产生，种族也必然过度繁衍；如果是这样，哪个天或地能容纳如此庞大且不断增长的众神群体？但我们该如何谈论那些将天神描绘为与姐妹女神乱伦结合，并指控他们通奸和不洁的人？我们进一步满怀信心地宣告，这些神灵从人那里得到的尊荣和崇拜，伴随着放荡和荒淫的行为。再说一次，经验丰富且技艺高超的雕刻家，形成其设计的构想后，按照艺术规则完善其作品；但不久之后，仿佛忘记了自己，崇拜自己的创造物，将其当作不朽的神明来敬拜，尽管他承认自己，即这形象的作者和制造者，是一个会死的人。不仅如此，他们甚至展示那些他们认为不朽者的坟墓和纪念碑，将神圣的尊荣归于死者；不知道真正有福和不朽的存在不需要可朽之人所能给予的区分：因为那只能被心灵之眼看见、只能被理智体悟的存在，无需藉外在形式来区分，也不容许任何形象来代表其特质和相似。但我们所说的尊荣是给予那些屈服于死亡权势的人：他们曾经是人，活着时是必死身体的寓居者。\par

% structure: p0010 source=h2 target=section reason=relative-heading-rank; base=h2

\section{第五章 \Emph{基督，圣子，创造了万有，并给每件事物指定了其存在的期限}}

然而，当我意在颂扬真天主的光荣时，为何还要以不洁的话语玷污我的舌头呢？倒不如让我藉着那从我赞颂之天主的永恒德能泉源涌出的纯净水流，洗净这苦涩的饮剂。我的特殊职责乃是借着生活中的行动，以及因祂所赐予的丰盛而显著的恩惠而应献上的感恩，来光荣基督。因此，我断言：祂奠定了这宇宙的根基，构思了人类族类，并以祂的圣言安排这一切。随即，祂将我们初造的父母（起初按祂的旨意，不识善恶）迁移到一个幸福的区域，那里繁花似锦，果实丰盈。最终，祂为他们指定了一个适于理性受造物的居所，然后向他们的理智——作为有灵之体——启示了善恶的知识。那时，祂也命令人类繁衍；世上每个健康的区域，直至环绕海洋的边界，都成了人的居所；而随着人口的增加，实用技艺的发明也随之并进。与此同时，各种较低等的动物也按适当比例增加，每一种类都显露出某种特性，即自然的特殊恩赐：驯良的以温顺和服从人类为特征；野性的以力量和敏捷为特征，还有一种本能的预见，警告它们逃避危险。祂将较温顺的动物完全置于人类的保护之下，但使人必须与那些较凶猛的动物争斗。接着，祂创造了有羽毛的族类，数量繁多，品性与习性各异；它们以各种颜色闪耀着光彩，并被赋予了天生的旋律能力。最后，祂以明智的判别安排了这世界范围内所含的一切其他事物，并为每个受造物指定了其存在的固定期限，从而完成了这完美整体的美丽秩序。\par

% structure: p0012 source=h2 target=section reason=relative-heading-rank; base=h2

\section{第六章　\Emph{关于命运的一般意见的虚妄性} \Emph{由人类法律及创造工程——其进程并非偶然，而是依照表明造物主设计的井然有序的安排——所证明。}}

然而，绝大多数人愚昧地将宇宙的秩序归诸自然，有些人则想象命运或偶然是原因。至于那些将万有归于命运的人，他们不知道用这个词时，他们只是吐出一个词语，并没有指明任何行动的力量，也没有指明任何真实而实在的存在。因为如果自然是万物的第一因，这命运本身又被当作什么？或者，如果命运的法则不可侵犯，我们该认为自然本身又是什么？的确，断言有命运的法则，就意味着这法则是立法者的作品：因此，如果命运本身就是法则，那它必是天主制定的法则。所以，万物都服从于天主，没有在他能力之外的东西。如果说命运是天主的旨意，并被这样认为，我们承认这一事实。但是，义德、节制或其他德行在何等程度上依赖于命运？如果这样，它们的对立面，如不义和放纵，又从何而来？因为恶源于自然，而非德行；德行是对自然性格和倾向的适当规范。但是，即使承认行为的结果——无论本身正确还是错误——取决于幸运或命运：正义的一般原则，即给每人应得的，又怎能归诸命运？或者说，法律、对德行的鼓励和对邪恶的劝戒、赞美、责备、惩罚，简言之，凡能激励德行、阻止恶习的，又怎能说源于幸运或偶然，而非源于正义——那是上智安排的天主的一个特有属性？因为发生在人身上的事件，是他们生活过程的结果。因此瘟疫或叛乱、饥荒和丰年交替出现，清楚而有力地表明这一切都是按照我们的生活方式调节的。因为天主喜欢善，厌恶一切不敬；他垂顾谦卑的人，却憎恶傲慢，以及那使受造物抬高自己的骄傲。虽然这些真理的证据在我们眼前清楚显明，但当我们收集并仿佛将思想集中于自身，深入思考其原因时，它们就显得更加清晰。我说，因此我们应当过着谦逊温和的生活，不让我们的思想骄傲地高过我们本性的状况，并时刻记得天主就在我们身边，他观察我们的一切行为。但我们还要进一步检验关于宇宙秩序取决于偶然或意外这个命题的真实性。难道我们要认为星星和其他天体、地和海、火和风、水和空气、季节的交替、夏冬的循环——这一切都是无目的和偶然的存在，而不是出于天主的创造之手吗？有些人甚至愚蠢到说这些东西大部分是人类因需要而发明的。即使我们承认这种说法在关于尘世和可朽坏的事物上似乎合理（虽然自然本身慷慨地提供一切美善），但我们能相信那不朽和不变的事物是人的发明吗？这些以及所有其他超越我们感官、只由理性领悟的事物，其存在并非来自人的物质生命，而是来自天主理智而永恒的本质。此外，这些事物的有序安排是他的上智安排之工：例如，日子从太阳获取光辉而明亮；夜晚在其落山后降临，而星宿的群队使夜晚本身免于完全黑暗。我们该怎样说月亮呢？它离太阳最远、相对时充满光明，而接近太阳时就亏损？这些岂不明显显明天主的理智和精明的智慧吗？再加上太阳光线必要的温暖，使地上的果实成熟；风的气流对季节的丰饶大有裨益；凉爽的阵雨；以及这一切事物的和谐，按照这和谐，一切都被合理而有条理地进行：最后，行星的永恒秩序，它们在指定时间返回同一位置：所有这些，以及星辰的完美服务，都服从于神圣的法则，难道不都是天主规定的明显证据吗？再者，山峰、幽深的谷地、平坦广阔的平原，既有用又悦目，难道能说它们独立于天主的旨意而存在吗？或者，土地和水的比例与交替，一者用于农耕，另一者用于运输我们所需的外来产品，岂不清晰证明了他精确而合比例的上智安排？例如，山上储水，平地接收，在吸收了足够更新土壤的水分后，将剩余的水送入海，海又将其递送到大洋。而我们竟敢说这一切都是偶然和意外发生的，虽然我们不能指出这偶然用什么形状或形式来标识——这东西既无理智也无感官的存在，只是作为空洞名称的声音在我们耳边回响！\par

% structure: p0014 source=h2 target=section reason=relative-heading-rank; base=h2

\section{第七章。\Emph{对于超越我们理解的事物，我们应光荣造物主的智慧，并将其原因唯独归于他，而非归于偶然。}}

事实上，\Quote{\OriginalTerm{偶然}{chance}}一词表达了那些以随意和非逻辑方式思考之人的说法；他们无法理解这些事物的原因，并由于自身领悟力的薄弱，而认为那些他们无法说明理由的事物是没有理由地被安排的。无疑地，有些事物具有奇妙的自然属性，其充分领悟十分困难：例如，温泉的性质。因为没有人能轻易解释如此强大之火的原因；而且确实令人惊奇的是，尽管四周被冷水包围，它却丝毫未失其天然热量。这些现象似乎在世界各地罕见发生，我相信，其目的是为人类提供令人信服的证据，证明那安排热与冷这两种直接相反的本性从同一源头流出之眷顾的权能。确实，许多，甚至数不胜数，是天主为了人的安慰和享受而赐予的恩赐；其中橄榄树和葡萄树的果实值得特别关注：一者因其更新和振奋灵魂的能力，另一者因其服务于我们的享受，也适用于治疗身体疾病。同样奇妙的，是河川的奔流，日夜不停地流动，呈现出一种永恒流动、永不停止的生命类型；白昼与黑夜的交替也同样奇妙。\par

% structure: p0016 source=h2 target=section reason=relative-heading-rank; base=h2

\section{第八章 \Emph{论天主慷慨地供应一切适合人之所需，却吝于满足其享乐；两样皆以人之利益为念。}}

以上所述足以证明，没有事物是无理由和无理智地存在的，而且理性本身和眷顾属于天主。祂也按适当比例分配了金属，如金、银、铜及其他；慷慨地供应那些最需要且普遍使用的，同时以宽厚而节俭的手分施那些仅仅服务于享乐和奢华装饰的，在吝啬与浪费之间保持中庸。因为，假若用于装饰的金属像其他金属一样容易获得，寻矿者就会因贪婪而轻视并忽略收集那些如铁或铜等适用于农耕、建房或装备船只的金属，而只关心那些导致奢侈和过度财富的金属。故此，如人们所说，寻找金和银比寻找其他金属远为困难和艰苦，辛劳的剧烈正好平衡了欲望的剧烈。还能再列举出多少天主眷顾的作为呢？它在所有如此慷慨赐予我们的恩赐中，明确地敦促我们实践节制和其他一切美德，并引导我们远离不当的贪求！要追溯这一切事物的隐秘原因，确实超出了人类理智的能力。因为一个脆弱而短暂之物的理智，如何能认识完整的真理，或纯粹地理解天主从起初的旨意呢？\par

% structure: p0018 source=h2 target=section reason=relative-heading-rank; base=h2

\section{第九章 \Emph{论哲学家们因渴望无所不知而陷入错误观念，其中某些人甚至陷入危险。—— 亦论柏拉图的学说。}}

因此，我们应当追求我们能力范围内且不超出本性限度的事物。因为辩论的说服力往往会将我们大多数人从事物的真理中引开，许多从事推理和自然科学研究的人就遇到这种情况，每当事物的宏大超出他们探究的能力时，他们就采用各种手段来掩盖真理。因此，他们的判断各异，彼此学说争辩对立，尽管他们自诩智慧。由此也引发了民众的骚动，当权者担心世袭制度颠覆而作出的严厉判决，对许多论辩者本身造成了毁灭。例如，\Person{苏格拉底}，因善辩而得意洋洋，沉溺于把坏的说成好的，不断玩弄争论的诡辩，结果成为同国同胞诽谤的牺牲品。\Person{毕达哥拉斯}，尤其声称拥有沉默和自律的美德，却被证明说了假话。因为他向意大利人宣称，他在埃及旅行期间接受的教义（很久以前已由该国的祭司们透露）是来自天主的个人启示。最后，\Person{柏拉图}本人，最温和优雅的一位，他首先尝试将人的思想从可感事物引向理智和永恒的对象，教导他们渴慕更崇高的思考，他最初真实地宣称有一位超越一切本质的天主，但后来又加上了第二位，在数目上区分为二，尽管两者拥有同一完美，且第二位神的存在出自第一位。因为他是宇宙的创造者和掌管者，显然是至高的；而第二位作为他命令的顺从执行者，将一切受造的起源归于他作为原因。因此，按照最正确的理性，我们可以说有一位眷顾和照管万有的存在，就是天主圣言，祂安排了万物；但圣言既然是天主本身，也就是天主子。因为除了\Quote{子}这一称号，我们还能用什么名称来称呼祂而不落入最严重的错误呢？因为圣父理所当然地被视为祂自己的圣言之父。至此，\Person{柏拉图}的观点是正确的；但随后他似乎偏离了真理，因为他引入了众多神祇，并为每位指定了具体形式。这给人类中愚昧的部分提供了更大的错误机会，他们不顾至高天主的眷顾，却崇拜自己设计、按照人或其他生物的样式制造的偶像。由此可见，这位哲学家的超越本性和令人钦佩的学识，尽管染上了这些错误，却绝非纯净无杂。然而，在我看来，他似乎又撤回并修正了自己的话，当他明确宣称理性的灵魂是天主的嘘气，并将万物分为两类：理智的和可感的：一方面单纯，另一方面由身体结构组成；一类仅凭理智理解，另一类凭判断和感觉估量。因此，前者分享了神圣精神，是单纯和无形体的，是永恒的，承受永生；而后者完全分解为构成它的元素，在永生中无份。他还教导令人钦佩的道理：那些度过美德生活的人，即善良圣洁之人的灵魂，在脱离身体后，被安置在天上最美丽的居所。这道理不仅令人赞叹，而且有益。因为谁能相信这样的陈述并渴慕如此幸福，却不愿实践正义和节制，远离恶习呢？与这道理一致，他描绘恶人的灵魂像漂浮的残骸一样在\Place{阿刻戎}和\Place{皮里弗勒格通}的河流上颠簸。\par

% structure: p0020 source=h2 target=section reason=relative-heading-rank; base=h2

\section{第十章。\Emph{论那些拒绝哲学家的学说以及圣经的学说者：以及我们应当相信诗人们的一切，还是全不信。}}

然而，有些人是如此愚昧，以至于当他们遇到这样的见解时，既不回心转意，也不感到惊惶；反而轻蔑嘲笑，仿佛在倾听寓言的虚构；也许称赞其辞藻的优美，却憎恶其训诫的严厉。然而，他们却相信诗人的虚构，使文明和野蛮之地都响彻已被驳斥的虚假故事。因为诗人宣称，死后灵魂的审判被委托给那些人——他们声称其出身源于诸神，赞美他们的公义与公正，并将他们描绘为死者的守护者。这些诗人还描述诸神之间的战斗及他们之间的某些战争习俗，并说他们受制于命运的力量。其中一些神祇被描绘为残忍的，另一些则对人类漠不关心，还有些具有可憎的品性。他们也描述诸神哀悼自己子女的杀戮，从而暗示他们无力拯救——不仅是陌生人，就连至亲也是如此。他们还描述他们受人类激情的支配，歌颂他们的战争与创伤、欢乐与悲伤。所有这些，他们似乎都值得相信。因为如果我们假设他们被神圣的冲动所驱使而尝试诗歌艺术，我们就必须相信他们，并被他们在这种灵感下所说的话说服。那么，他们谈论的是他们的神祇所遭受的灾难；这些灾难当然完全是真实的！但有人会反驳说，撒谎是诗人的特权，因为诗歌的独特领域是取悦听众的心灵，而真理的本质在于所讲述的事情必须与实际完全一致。让我们承认，诗歌偶尔有隐藏真相的特性。但那些说谎者并非没有目的；他们要么出于对个人利益或好处的渴望，要么可能由于意识到某些邪恶行为，而因害怕法律报复的威胁而被迫掩盖真相。但依我之见，他们完全有可能——至少在论及至高者的本质时忠实于真理——从而避免说谎和不敬的罪责。\par

% structure: p0022 source=h2 target=section reason=relative-heading-rank; base=h2

\section{第十一章 \Emph{论我们主的降生成人：其性质与原因}}

凡行为与美德生命不相称，自知生活纷乱无序者，让他悔改，以开启的灵性目光转向天主；让他放弃过去的恶行，即使到晚年获得智慧也已满足。然而，我们并未从人的教导得到帮助；不，凡有识者所称誉的品性恩宠，全是天主的恩赐。我能用来对抗\Person{撒殚}军械库中致命武器的盾牌并不薄弱，我指的是我所拥有的那令祂喜悦的知识：我将从中选取适合我当前意图的，继续歌颂万有之父。但是，人类救主基督，请临在帮助我完成这神圣任务！引导颂扬你德行的言语，教导我相称地称扬你的赞美。现在，不要期望听到优雅言辞的恩典：因为我深知，那些只为悦耳而说话、追求喝彩而非稳健论证的人的软弱修辞，对明辨的听众来说是令人厌恶的。因此，有些亵渎而愚蠢的人断言，我们所敬拜的基督被公正地判处死刑，而那位为万有生命之源者，他自己被剥夺了生命。那些曾经敢于踏上不敬之路、抛弃一切畏惧、毫不掩饰自身堕落的人做出这样的断言，并不令人惊讶。但这超越愚蠢本身的界限，他们似乎真能说服自己：不朽坏的天主屈服于人的暴力，而非屈服于祂对人类所怀的独一之爱；他们未能察觉天主的宽厚和忍耐不因任何侮辱而改变，不因任何咒骂而偏离其固有的坚定；而是始终如一，以智慧之灵和灵魂的伟大，击溃和击退攻击者的野蛮凶残。天主的仁慈恩惠已决意废除不义，高扬秩序与正义。因此，祂聚集了人间最智慧的一群，制定了那最高贵、最有益的教义，旨在引导人类的善良与蒙福者效法祂自己的眷顾。我们能说出比这更高的祝福吗：天主规定正义之路，使那些被认为配得祂教导的人肖似祂自己；以致良善可传于各阶层人类，而永恒之福为其结果？这是光荣的胜利！这是真正的权能！这是伟大的工程，与其作者相称，即恢复所有人健全的思想：我们欢喜地将这胜利的荣耀归于你，万有的救主！但是你呢，卑劣可恶的亵渎，你的荣耀在于谎言、谣言和诽谤；你的能力是欺骗和战胜缺乏经验的青年，以及那些仍保留青年愚昧的人。你引诱他们离开对真天主的服事，设立假偶像作为他们敬拜和祈祷的对象；因此他们愚昧的报应等待着你的迷惑受害者：因为他们诽谤基督——一切祝福之源，祂是天主，也是天主子。世上最优秀、最智慧的民族之敬拜，难道不是正当地指向那位天主吗？祂拥有无限权能，却坚定不移地忠实于自己的旨意，并保持其特有的慈爱与对人的爱毫不减损？那么，你们这些不敬神的人，趁着对你们过犯的报复尚未施行，你们仍然可以离开；去你们的祭献、宴席、狂欢和醉酒场所吧，在那里，在宗教的外表下，你们的心专注于放荡享乐，假装献祭，而你们却是自己情欲的甘心奴隶。你们对任何善都没有认识，甚至对大能天主的首要诫命也一无所知，祂向人宣告祂的旨意，并委托祂的圣子指引人类生活的道路，以便那些经过德行和自制生涯的人，按照那圣子的审判，可以获得第二次生命，甚至是有福和幸福的存在。我现在已经宣告了天主关于祂为人类规定之生命的旨令，既不无知（像许多人那样），也不基于意见或猜测。但或许有人会问：这圣子的称号从何而来？我们所说的这发生从何而来，如果天主确实是唯一且不能与另一位结合？然而，我们应该将发生视为两类：一类是自然生育的方式，人所共知；另一类是永恒原因的效果，其方式由天主的前知和祂所爱的人中的那些人所见。因为智慧的人会认识到调节创造和谐的原因。因此，既然没有事物没有原因而存在，存在实体的原因必然先于它们的存在。但既然世界及其包含的一切存在并被保存，它们的保存者必先存在；因此基督是保存的原因，万物的保存是效果：正如圣父是圣子的原因，圣子是那原因的效果。那么，已经说得足够证明祂存在的优先性。但我们如何解释祂降世到人间呢？祂这样做的动机，正如先知所预言的，源于祂对所有利益的细心关怀：因为创造者必须关怀祂自己的作品。但当时候到了，祂需要取一个尘世的身体，并寄居于此世，由于需要，祂为自己设计了一种新的诞生方式。受孕存在，却没有婚姻；分娩，却是纯洁的童贞：一位少女成了天主之母！永恒的本性接受了时间存在的开始：一种精神本质的可感知形式，非形体光明的物质显现，出现了。伴随这一伟大事件的情境同样奇妙。一只光辉的鸽子，如同从\Person{诺厄}方舟飞出的那只，落在\Person{童贞玛利亚}的胸前：与这不可捉摸的结合相一致，其结果比贞洁更纯洁，比天真更无邪。从婴幼时期就拥有天主的智慧，在\Place{约旦河}受洗时被以敬畏之心接纳，蒙受了那王室傅油，即普遍智能之神；拥有知识和能力行奇迹，医治超越人类医术的疾病；祂迅速无碍地应允人的祈祷，祂的一生确实毫无保留地为人类的福祉奉献。祂的教义灌输的不仅是谨慎，而是真正的智慧：祂的听众被教导的不仅是纯粹的社会美德，而是通往灵性世界的道路；他们致力于默观不变和永恒的事物，以及认识至高圣父。祂赐予的恩惠不是普通的祝福：给失明者视觉的恩赐；给无助的软弱者健康的活力；代替死亡，重获生命。我不详述在旷野中那丰富的供应，藉此少量的食物成为一大群人所需要之充足而持久的供应。因此，我们按照自己微弱的能力向您，我们的天主和救主，献上感谢；向您，基督，全能圣父的最高眷顾，您既拯救我们脱离邪恶，又将您最蒙福的教义传授给我们：因为我说这些，不是为了赞美，而是为了感恩。因为哪个凡人能相称地宣告对你的赞美？我们从你得知，你从无中召叫了创造，并以你的光明照亮了它；你以和谐与秩序的法则规范了元素的混沌。但我们尤其注意到你的慈爱，因为你使那些心倾向你的人热切渴望神圣而有福的生命，并提供了他们如同真正祝福的商人，将所获得的智慧和好运传递许多其他人；同时，他们自己收获美德的永恒果实。从恶习的束缚中解脱，充满对同胞的爱，他们常将慈悲置于眼前，盼望信德的应许；致力于谦逊，以及人类过往生涯所抛弃的一切美德[但如今由那位眷顾万有的祂所恢复]。找不到其他力量为如此邪恶、以及此前宣称统治人类的不义之灵设计补救办法。然而，眷顾甚至能到达此处的境况，并轻易恢复了那些被暴力和人类情欲的放纵所扰乱的一切。祂毫无掩饰地行使了这种恢复能力。因为祂知道，虽然有些人能认出并理解祂的权能，但另一些人其兽性和无感觉的本性使他们完全依赖自己感官的见证。因此，在光天化日之下，为使无论善恶，无人有怀疑的余地，祂显示了自己有福而奇妙的治愈权能：使死者复活，并以一句话恢复那些丧失身体感觉者的能力。简而言之，我们能否认为，使海变得坚固如实地，平息风暴的狂怒，最终在通过行这些奇妙作为将人的不信转为坚定信德之后升天，所要求的少于全能权能，少于天主的工作？祂受难之时也伴随同样奇事：当太阳变暗，黑夜的阴影遮蔽了白昼之光。那时恐惧抓住了各处的人民，认为万物的结局已经来到，像创造秩序开始之前那样的混沌将再次得势。当时，人们也寻找如此可怕灾祸的原因，以及人的过犯在哪些方面激起了上天的愤怒；直到天主自己，祂以平静的尊严审视不敬神者的傲慢，更新了天空的面貌，以群星装饰了它。这样，阴云密布的自然面容再次恢复了她原来的美丽。\par

% structure: p0024 source=h2 target=section reason=relative-heading-rank; base=h2

\section{Chapter 12. \Emph{对奥秘无知者；其无知出自自愿。认识它的人所等待的福分，尤其是那些在宣认信仰中死去的人。}}

但有些爱亵渎的人会说，天主本有能力改善和柔化人的自然意志。我问，有什么比与天主亲自交谈更好的方法、更有效的努力来挽回恶人呢？他不是明明在场教导他们德行的原则吗？如果天主亲身的教导都没有效果，那么他若持续缺席、不被听闻，情况岂不更糟？那么，是什么有能力阻碍这至福的教义？是人的乖谬愚妄。因为每当我们以愤怒的不耐烦去接受那些为我们的福祉和益处而赐下的训诫时，我们理解的清晰度就立刻变得模糊。事实上，人正是选择了忽视这些训诫，对那令他们厌恶的诫命充耳不闻；然而，如果他们听了，就会获得一个值得如此关注的赏报，而且不仅是为了今生，更是为了来世——那才是唯一真实的生命。因为顺从天主的赏报是不朽的永生，凡认识他的人都能渴望得到，且他们安排自己的生活，以便为他人树立榜样，并作为那些渴望在德行上卓越之人的永恒效法标准。因此，这教义托付给了有智慧的人，使他们所传讲的真理能被家庭成员以谨慎和纯洁的良心保存，从而确保对天主诫命真实而坚定的遵守，其果实就是那种因纯正的信德和在天主面前真诚的圣洁而产生的面对死亡的勇气。如此武装的人能抵御世界的风暴，甚至被天主无敌的能力支撑到殉道，借此他勇敢地克服最大的恐惧，并被那位他以崇高方式见证的主视为配得上荣耀的冠冕。他自己也不居功，深知正是天主赐予他力量，使他能忍受并以热忱完成神圣的命令。这样的行为理应获得永不磨灭的纪念和永远的光荣。因为殉道者的生活是清醒和顺从天主旨意的，他的死亡则是真正伟大和慷慨灵魂坚毅的榜样。因此，随之而来的是赞美全视天主的赞歌、圣咏、话语和颂歌；并为纪念这样的人献上感恩的祭献，一种不流血、无害的祭献，其中不需要芬芳的乳香，不需要火；只需足够纯净的光来满足聚集的朝拜者。还有许多具有爱德精神的人，他们为救济穷人、援助被赶出家园的人而预备节制的宴席：这种习俗只有那些心思不与神圣的教义相符的人才会视为负担。\par

% structure: p0026 source=h2 target=section reason=relative-heading-rank; base=h2

\section{Chapter 13. \Emph{受造物之间存在必然差异。向善与向恶的倾向取决于人的意志；因此，审判是必要而合理的事。}}

的确，有些人竟以幼稚的傲慢敢就此事指责天主，问为什么他不为所有人创造完全相同的自然性情，却反而命定许多不同甚至相反性质的事物存在，从而造成我们道德行为与性格的差异。难道（他们说）如果所有人都具有相同的道德品性，不是更有利于顺从天主的命令、正确领会他本身、以及坚固个人的信德吗？指望这种情况发生并忘记了世界的构造与世界内事物的构造不同、物理对象与道德对象性质不同、身体的情感与灵魂的情感不同，实在是荒谬的。[因为不朽的灵魂在尊贵上远超物质世界，比起可朽坏和尘世的受造物更为有福，按其与天主的相似和接近程度。] 人类也并非被排除在天主的美善之外；不过这不是所有人不加区别的份，而只属于那些深入探索天主本性、并以认识神圣事物作为生活首要目标的人。\par

% structure: p0028 source=h2 target=section reason=relative-heading-rank; base=h2

\section{Chapter 14. \Emph{受造的本性与非受造的存在者无限不同；人靠着德行生活能最接近非受造者。}}

将受造之物与永恒之物相比，确实是极端的愚昧，因为永恒之物既无开始也无终结，而受造之物既被创造、被召唤而存在，并获得其存在在某个特定时间的开始，那么它们必然也必须有一个终结。这些被造之物怎能与那命令它们存在的天主相比呢？若是这样，他便不能被正确地视为具有命令它们存在的能力了。天上的事物也无法与他相比，正如物质世界无法与理智世界相比，或复制品无法与原型相比一样。不仅如此，将万物如此混淆，借将他与人类甚至禽兽相比来遮蔽天主的尊荣，岂不荒谬？而渴望拥有与天主同等的权能，岂不是那些完全脱离了节制与道德生活的疯人的特征吗？若我们在某种程度上的确渴望如天主般的福乐，我们的职责便是依照他的诫命生活：这样，在完成符合他规定的法律的生活之后，我们就能永远住在永恒不朽的居所中，超脱命运的力量。因为人身上唯一能与天主媲美的能力，是真诚无伪的服务和全心奉献，以及默观和研习一切令他喜悦之事，将我们的情感提升到超越尘世之物，并尽可能将我们的思想指向崇高和天上的对象：因为据称，从这样的努力中，我们获得了一场比许多祝福更宝贵的胜利。因此，受造之物在尊严和能力上不平等的原因，正如我所描述的。明智之人对此满怀感恩和喜乐而顺从；而那些不满之人则显露出自己的愚昧，他们的傲慢将得到应有的报应。\par

% structure: p0030 source=h2 target=section reason=relative-heading-rank; base=h2

\section{第十五章。\Emph{论救主的教义与奇迹；以及他赐予那些服从他的人的恩惠。}}

天主圣子邀请众人从事德行的实践，并将自己呈现给所有有悟性的人，作为他救恩训诫的导师。除非我们确实自欺，并悲惨地不知这一事实：为了我们的益处，即为了确保人类蒙受祝福，他巡游世间；并将当世最优秀的人召聚在周围，将充满利益且有能力保守他们行在德行生活之途的训诲托付给他们；教导他们信德与义德，这些是对抗那恶神敌对势力的真正良药，那恶神乐于诱捕并欺骗无经验者。因此，他探望病患，解除虚弱者的痛苦，安慰那些陷于极端贫困与匮乏的人。他也称赞健全理性的节制品格，命令他的追随者以尊严和耐心忍受各种伤害与蔑视：教导他们视之为圣父所允许的考验，而胜利常属于那些高贵承受降临于己的灾祸的人。因为他向他们保证，最高的力量全在于灵魂的坚毅，结合那无非是对真理与善之认知的哲学，在人心内产生慷慨的习惯，将他们以正当方式获得的财富与较贫穷的兄弟共享。同时他彻底禁止一切傲慢的压迫，宣告：正如他来与卑微者交往，那些轻视卑微者的人将被排除在他的恩宠之外。这就是他用以考验那些服从他权柄者的信德的试炼，如此他不只使他们准备好藐视危险与恐怖，同时也教导他们对他有最真诚的信赖。有一次他也出言责备，以抑止一位同伴的热心，那人过于轻易屈从于冲动的激情，持剑攻击，急切想要保护救主性命，却暴露了自己。那时他命他住手，并将剑收回鞘中，指责他不信任在他内的庇护与安全，并庄严宣告凡试图以同样的攻击报复伤害、或用剑者，必死于暴亡。这的确是属天的智慧，宁愿忍受而非施加伤害，并准备好在必要时受苦，而非作恶。因为伤害行为本身就是最严重的恶，承受最重惩罚的不是受害方，而是加害方。的确，一个顺从天主旨意的人，若能坚定信赖那常以援助保护其仆人免受伤害的天主的保护，就能避免施害与受害双重的恶。因为那信赖天主的人，怎么会试图从自身寻求资源呢？在这种情况下，他必须忍受胜负未卜的战斗；任何有智慧的人都不会选择不确定的结果而非确定的结果。再者，那经历过多种危险，且时常因天主轻轻点头就轻易脱离一切恐惧的人，那仿佛穿越了由救主之言所铺平、为百姓提供坚固道路的海的人，怎么会怀疑天主的临在与援助呢？我相信，这就是信德的可靠基础、信心的真正根基：我们发现这样的奇迹，是在上智安排的天主的命令下完成并成就的。因此，即使在考验中，我们也找不到理由后悔自己的信德，而是在天主内保持不摇动的望德；当这种信心的习惯在灵魂内建立时，天主自己便住在最深的思念中。但他是不可战胜的：因此，灵魂内若有这样一位不可战胜者，就不会被周围可能有的危险所克服。同样，我们从天主自己的胜利学到这真理：他一心为人类谋求祝福，虽遭不虔之人的恶毒严重侮辱，却毫发无伤地经历了他的苦难，并获得巨大的胜利、不朽的凯旋冠冕，战胜一切不义；如此，他实现了自己上智安排与爱德对义人的目的，并摧毁了不虔与不义者的残忍。\par

% structure: p0032 source=h2 target=section reason=relative-heading-rank; base=h2

\section{第十六章. \Emph{先知们预言了基督的来临；并注定要推翻偶像和崇拜偶像的城市。}}

很早以前，他的苦难以及他在肉身中的来临，就已被先知们预言了。他降生成人的时间也曾被预言，以及毁灭不义与放荡的果实（这些果实对义德的行为和道路具有毁灭性）的方式，并让整个世界分享智慧与健全明辨的德性，通过救主要颁布的那些行为准则在人的心意中几乎普遍的盛行；藉此，对天主的敬拜得以确认，而迷信的礼仪被彻底废除。通过这些，不仅设计了宰杀牲畜，还设计了人祭和受诅咒的敬拜的污秽：例如，依照亚述和埃及的法律，无辜者的生命被献在铜或土制的偶像上。因此，这些民族受到了配得上如此污秽敬拜的报应。\Place{孟斐斯}和\Place{巴比伦}据宣告将被荒废，与它们祖先的神明一同沦为废墟。我这样说，不是听他人报告，而是我亲自在场，实际看见了这些城市中最悲惨的、不幸的孟斐斯。\Person{梅瑟}遵照天主的命令，摧毁了昔日强大的\Person{法郎}的土地（法郎的傲慢导致了他的毁灭），并摧毁了他的军队（这军队曾战胜众多强大的民族，是一支防御坚固、武器精良的军队），不是靠箭矢的飞射或敌对武器的投掷，而只靠圣洁的祈祷和安静的恳求。\par

% structure: p0034 source=h2 target=section reason=relative-heading-rank; base=h2

\section{第十七章. \Emph{论梅瑟的智慧，它是异教民族中的智者所效法的对象。也论达尼尔与三个青年。}}

没有任何民族曾像\Person{梅瑟}所领导的那一民族一样，受到如此高深的祝福；若不是他们自愿远离圣神的引导，他们本会继续享受更高的祝福。然而，谁能恰当地歌颂\Person{梅瑟}本人的赞美？他治理了一个不驯的民族，陶冶他们的心灵，使之养成服从和尊敬的习惯，将他们从被俘中解救出来，恢复自由，将他们的哀伤化为喜乐，并使他们的心灵如此提升，以至于由于与前情况的极大反差和富足的繁荣，民众的精神变得傲慢和自负。他在智慧上如此超越前人，以至于连异教民族所赞扬的智者和哲学家都渴望效法他的智慧。因为\Person{毕达哥拉斯}效法他的智慧，达到了如此的自制程度，以至于他成为\Person{柏拉图}——本身就是审慎的典范——自制力的标准。再者，那位古\Place{叙利亚}王的残忍是何等巨大可怕，\Person{达尼尔}战胜了他，这位先知揭示了未来的奥秘，他的行为彰显了超凡的心灵伟大，他品格和生活的光辉超越众人。这位暴君的名字是\Person{拿步高}，他的后代后来灭绝了，他那庞大而强大的权力转移到了\Place{波斯}人手中。这位暴君的财富在当时甚至现在都远近闻名，还有他不合时宜地沉迷于非法崇拜、他的偶像雕像，那些雕像高耸入云，由各种金属制成，以及为维护这种崇拜而制定的可怕而野蛮的法律。\Person{达尼尔}，因对真天主的真诚虔诚而得到支持，完全鄙视这些恐怖，并预言暴君不合时宜的热心将对他自己产生可怕的恶果。然而，他未能说服暴君（因为过度的财富是真正判断健全的有效障碍），最终，君主展示了其性格的野蛮残忍，命令将这位正义的先知暴露在野兽的狂怒之下。确实，那些兄弟（后来他人效法他们的榜样，因对救主之名的信德而获得超凡荣耀）所表现出的团结精神也是崇高的，我指的是那些在火窑中安然无恙的人，以及那些被指定吞噬他们的恐怖，他们用身体的神圣接触驱散了围绕他们的火焰。在\Place{亚述帝国}被天上雷霆摧毁后，天主的眷顾将\Person{达尼尔}带到\Place{波斯}王\Person{冈比西斯}的宫廷。然而，嫉妒甚至在这里跟随了他；不仅是嫉妒，还有术士对他性命的致命阴谋，以及一连串许多迫切的危险，所有这些他都因基督的眷顾照顾轻易地脱身，并在各种德行的实践中闪耀光辉。他每天三次向天主献上祈祷，并展示了他所显的超自然力量的可铭记的证据：因此，术士们对他祈祷的功效充满嫉妒，向国王陈述拥有这种力量的危险，并说服他将这位\Place{波斯}人民的杰出恩人判给野蛮的狮子吞噬。于是，\Person{达尼尔}这样被判，被投入狮穴（不是为了遭受死亡，而是为了赢得不朽的荣耀）；尽管被这些凶猛的野兽包围，他发现它们比将他关在里面的人更温和。因平静而坚定的祈祷的力量支持，他得以制服所有这些按其本性凶猛的动物。\Person{冈比西斯}得知此事后（因为如此强大的神圣力量的证据不可能被隐藏），对这件奇妙的故事感到惊讶，并对轻信术士的诽谤指控感到后悔，然而决定亲自见证这一景象。但当他看到先知高举双手向基督献上赞美，而狮子蹲伏在他脚下，仿佛在朝拜时，他立即判那些他所听信的术士遭受同样的判决，并将他们关进狮穴。那些曾经如此温顺的野兽立刻扑向他们的受害者，以它们全部的本性凶猛撕碎并毁灭了他们全部。\par

% structure: p0036 source=h2 target=section reason=relative-heading-rank; base=h2

\section{第18章 \Emph{论埃律特赖亚女先知，她在预言性的藏头诗中指向了我们的主及其受难。藏头诗是：\Quote{耶稣基督，天主之子，救主，十字架。}}}

然而，我甚至希望从外邦来源中获取一份关于基督神性的见证。因为在这见证上，显然连那些亵渎他名号的人也必须承认他是天主，是天主子——只要他们确实相信那些与自己意见一致的人的话。\Person{埃律特赖亚女先知}本人向我们保证，她生活在洪水后的第六代，是\Person{阿波罗}的女祭司，佩戴着神圣的束发带以模仿她所侍奉的神，还守护着被蛇缠住的三脚架，并对来到她神龛前的人给出预言性的回答；她被父母愚昧地奉献给这种服务，一种不能产生任何善良或高尚，只会产生不当狂乱的服务，正如我们在\Person{达佛涅}的故事中所看到的那样。有一次，她冲进她虚妄迷信的圣所，真的从上面得到了灵感，并以预言性的诗句宣告了天主的未来旨意；通过这些诗句的起始字母清楚地表明了耶稣的降临，形成了一首藏头诗，其词是：\Emph{耶稣基督，天主之子，救主，十字架}。这些诗句如下：\par

\Quote{审判！大地渗出的孔隙将标志那日；\newline{}高天的君王将彰显他的荣光：\newline{}万有的主宰，高踞于他的宝座，\newline{}无数群众将承认他们的天主；\newline{}将看见他们的审判者，怀着喜惧交织之情，\newline{}与诸圣同戴冠冕，以人的形体显现。\newline{}当大地的荣华衰败时，\newline{}财富、炫耀和人的偶像崇拜何其虚妄！\newline{}在那可畏的时刻，当自然的烈火命运\newline{}惊扰坟墓中沉睡的住户时，\newline{}凡血肉都将战栗站立；每一隐秘的诡计、\newline{}久被遗忘的罪、愧疚与欺诈的意念，\newline{}都将暴露在天主搜索的光明之下：\newline{}那时毫无避难所，只有绝望的苦痛。\newline{}黑夜的阴影将铺满天穹，\newline{}从大地、太阳、星辰和月亮夺去它们的光；\newline{}天主的臂膀将粉碎群山骄傲的巍峨；\newline{}海洋的平原不再有舰队航行。\newline{}水源枯竭，不再有河流的奔腾声\newline{}抚慰，没有泉源滋润焦干的土地。\newline{}周围、远方，号角声将滚动，\newline{}那是久被延迟的怒涛之声，终于显明。\newline{}在哑然的敬畏中，当大地根基呻吟时，\newline{}在审判座上，地上的君王将承认他们的天主。\newline{}那时，被高举，具有神圣威严，\newline{}光辉灿烂，看哪，救恩的标记！\newline{}那一位主的十字架，他曾为罪人舍己，\newline{}被人唾弃，如今被天地承认，\newline{}用他的铁杖统治万邦。\newline{}这就是这些神秘诗句所显示的名号；\newline{}救主，永恒的君王，除去了我们的罪。}\par

显然，这位贞女是在神圣感召下说出这些诗句的。我不得不认为她是有福的，因为救主如此拣选她来启示他对我们恩慈的旨意。\par

% structure: p0040 source=h2 target=section reason=relative-heading-rank; base=h2

\section{第十九章　《论关乎我们救主的这预言并非基督教会中某人的杜撰，而是埃里特莱亚女先知西比尔之见证，其书卷在基督降世前由西塞罗译成拉丁文。又论维吉尔也曾提及此事及贞女之子诞生；虽因惧怕当权者而将此奥秘隐约道出。》}

然而，许多承认\OriginalTerm{埃里特莱亚女先知西比尔}{Erythræan Sibyl}确是先知的人，仍拒绝相信这预言，以为某位信奉我们信仰且并非不谙诗艺的人，是这些诗句的作者。总之，他们认为这是伪作，被声称是西比尔的预言，因为其中含有有益的道德意旨，有助于约束放纵，引导人过清醒端正的生活。但在此事上，真理是明显的，因为我国同胞的勤勉已仔细计算了时间；因此没有理由怀疑此诗是在基督降临并被定罪之后所作，或普遍传言说这些诗句是西比尔在早年的预言是虚假的。因为人们公认\OriginalTerm{西塞罗}{Cicero}熟悉此诗，他将其译成拉丁文，并纳入自己的著作。这位作家在安东尼当权时被处死，安东尼又被奥古斯都击败，奥古斯都在位五十六年。提庇留继位，在他的时代，救主的降临照亮了世界，我们至圣宗教的奥秘开始盛行，可以说是新人类的开端。我想，拉丁诗人的宗主对此如此说：\par

\Quote{看哪，一个新的、天生的族类出现了。}\par

\Quote{西西里的缪斯，奏出更高昂的曲调。}\par

\Quote{库玛的神谕之声再次响起。}\par

显然指的是\OriginalTerm{库迈亚的西比尔}{Cumæan Sibyl}。甚至这还不够：诗人似乎不可抗拒地被迫作证，又进一步说。那么他说了什么呢？\par

\Quote{看哪！流转的年岁带来新的福佑：\newline{}贞女来了，与她同来的是久盼的君王。}\par

那么，那将要来的贞女是谁？岂不就是那位充满圣神并怀有圣神之子的吗？为什么怀有圣神之子的她不可能成为并永远保持童贞呢？这位君王也将再来，他的来临将减轻世界的忧伤。诗人又加说：\Quote{看哪！流转的年岁带来新的福佑：\newline{}贞女来了，与她同来的是久盼的君王。}\par

\Quote{你，纯洁的路喀娜，迎接这新生的婴孩，\newline{}在他的统治下，铁的世代终结，\newline{}黄金的后裔从天降临；\newline{}他的国度将恢复流放的德行，\newline{}罪恶不再威胁有罪的世界。}\par

我们看出这些话是明说同时又隐晦地以比喻方式说的。那些深入探索词句含义的人，能够辨别出基督的神性。但唯恐帝都中的任何有权势者可能指控诗人写了违背国家法律、颠覆自古流传的宗教信仰的东西，他有意地隐晦了真理。因为我相信，他知晓那赐予我们主以救主之名的有福奥秘；但为了躲避残酷之人的严厉，他将听者的思想引向他们所熟悉的事物，说必须为新生的婴孩竖立祭台、建造殿宇、献上牺牲。他的结束语也适应了习惯于这类信条的人的情感；因为他说：\Quote{你，纯洁的路喀娜，迎接这新生的婴孩，\newline{}在他的统治下，铁的世代终结，\newline{}黄金的后裔从天降临；\newline{}他的国度将恢复流放的德行，\newline{}罪恶不再威胁有罪的世界。}\par

% structure: p0050 source=h2 target=section reason=relative-heading-rank; base=h2

\section{第二十章　《再引维吉留斯·马罗论基督的语录及其解释，表明这奥秘是隐约指出的，正如诗人所预期的那样。》}

\Quote{他将过一种不朽的生命，并被英雄所见，\newline{}他自己也将看见英雄；}\par

显然指的是义人。\par

\Quote{纷争的万民他将在和平中约束，\newline{}并以父性的德行治理人类。\newline{}大地未经召唤就将结出最初的果实，\newline{}芬芳的香草，迎接她的君王婴孩。}\par

这位绝顶智慧和有教养的人的确熟知时代的残酷特点。他继续写道：\Quote{山羊，未经呼唤，将满载乳房回家；\newline{}低鸣的牛群不再惧怕猛狮。}\par

\Quote{山羊，未经呼唤，将满载乳房回家；\newline{}低鸣的牛群不再惧怕猛狮。}\par

说得好：因为信仰不会在皇宫中的权贵面前畏惧。\par

\Quote{他的摇篮将冠以盛开的鲜花：\newline{}蛇的后裔必死；圣地\newline{}将拒绝生长杂草和有毒植物；\newline{}每丛矮树都将佩戴亚述的玫瑰。}\par

没有什么比这更真实、更符合救主的卓越了。因为圣神的能力将天主的摇篮，如同芬芳的花朵，献给新生的族类。那蛇及其毒液也消亡了，它最初欺骗了我们的原祖父母，将他们的思想从本性的纯真引向享乐，好让他们尝到那威胁的死亡。因为在救主来临之前，蛇的能力显现在颠覆那些没有牢固希望、不知义人永生之人的灵魂上。但在救主受苦并暂时离开他所取的身体之后，复活的能力通过圣神的传递启示给人；人类罪恶的污秽，无论残留多少，都通过神圣洗礼的洗涤而被清除。\par

那时，救主确实能命令他的门徒放心，并记念他可敬而荣耀的复活，期待自己也同样复活。因此，毒蛇之族可以说已灭绝。死亡本身也灭绝了，复活的真理被封印。再者，亚述之族也消失了，它曾首先引导人走向对天主的信仰。但当他处处提到阿摩姆的生长时，他暗指真正敬拜天主者的众多。因为如同许多枝子，戴着芬芳的花朵，得到合宜的浇灌，从同一根上生出。最公正地说，\Person{玛罗}，你最智慧的诗人！而后面的一切都与之一致。\par

\Quote{但当英雄的德行使他倾听青春，\newline{}并学会尊敬父亲的德行。}\par

通过对英雄的赞美，他暗示义人的作为；通过他父亲的德行，他谈到世界的创造和永恒的结构；也许也谈到那些法律，天主可爱的教会被其引导，并在正义与德行的道路上被安排。再者，那介于善恶之间的生命状态向更高事物的进展是令人钦佩的，这种状态很少容许突然的改变：\par

\Quote{无需劳作的收成将装饰田野，}\par

就是说，神律的果实为人服务而生长。\par

\Quote{而串串葡萄将在每根荆棘上泛红。}\par

在人类生命腐化而无律的时期远非如此。\par

\Quote{多节的橡树将滴下蜜雨。}\par

他在这里描述了那个时代的人的愚昧和顽梗；也许他也暗示，那些为天主的原因受苦的人，将收获自己忍耐的甜美果实。\par

\Quote{然而，旧骗局的一些脚步将留存；\newline{}商人仍将为了利润耕犁深海；\newline{}大城市将被城墙围住，\newline{}锋利的犁刃将困扰肥沃的土地；\newline{}另一个\Person{提菲斯}将探索新的海洋；\newline{}另一艘\OriginalTerm{阿尔戈}{Argo}将把首领们送上\Place{伊比利亚}海岸；\newline{}另一个\Person{海伦}将产生新的战争，\newline{}伟大的\Person{阿喀琉斯}将催迫\OriginalTerm{特洛伊}{Trojan}的命运。}\par

说得对，最智慧的诗人！你恰恰将诗人的自由权用到了恰当的程度。因为你的目的不是承担先知的功能，你并无此权利。我想他也受到了威胁的危险的约束，那危险威胁着攻击古代宗教习俗信用的人。因此，他谨慎而尽可能安全地向那些能够理解的人呈现真理；当他斥责战争的防御工事和冲突（这些在人类生活中确实仍可找到）时，他描述我们的救主前往与\Place{特洛伊}作战，以\Place{特洛伊}指代世界本身。而他确实在与邪恶的反对势力抗争中坚持，那使命是由他自己眷顾的计划和他全能圣父的命令所差遣的。那么，诗人接下来的话是什么呢？\par

\Quote{但当他长成成熟的成年时，}\par

也就是说，当他到达成年，彻底消除包围人类生命道路的邪恶，并通过和平的祝福使世界安宁时：\par

\Quote{贪婪的水手将放弃海洋；\newline{}没有龙骨会为了外国商品划破波浪，\newline{}因为每一块土地都将生产每一种产品。\newline{}劳累的农夫将卸下他的牛；\newline{}没有犁会伤害田地，没有修枝刀会伤害葡萄树；\newline{}羊毛也不会在伪装的色彩中闪耀；\newline{}但羊群奢华的父，\newline{}以本色的紫和金，\newline{}在它壮丽的羊毛下将骄傲地出汗；\newline{}在提尔紫的袍子下羊羔将咩咩叫。\newline{}年岁成熟，向准备好的荣誉移动，\newline{}哦，属天之种，哦，\Person{朱庇特}的养子！\newline{}看，劳苦的大自然呼唤你支撑\newline{}天空、大地和海洋的摇动火焰！\newline{}看，恢复其基础，大地、海洋和空气；\newline{}欢乐的世代，从后面，在成群的行列中出现。\newline{}要赞美你，愿上天延长我的呼吸，\newline{}注入配得上如此歌谣的灵，\newline{}色雷斯的\Person{俄耳甫斯}也不能超越我的诗，\newline{}也非\Person{利诺斯}，戴着永不凋谢的桂冠；\newline{}虽然每人都应激发他的天上父母；\newline{}\Person{缪斯}教导声音，\Person{福玻斯}调谐竖琴。\newline{}如果\Person{潘}以诗句竞争，而你是我的主题，\newline{}\Place{阿卡迪亚}的裁判将判他们的神有罪。}\par

\Quote{看（他说），强大的世界和元素如何一同表示喜乐。}\par

% structure: p0074 source=h2 target=section reason=relative-heading-rank; base=h2

\section{第二十一章。\Emph{这些事不能是对一个纯粹的人说的；而不信者，由于对宗教的无知，甚至不知道自己存在的起源。}}

有人可能会愚昧地以为这些话是对一个普通凡人的出生说的。但如果仅此而已，有什么理由大地不需要种子或犁，葡萄树不需要修枝刀或其他栽培手段？我们怎能设想这些是对一个凡人的出生说的？因为大自然是神圣旨意的仆役，不是服从人的命令的工具。确实，元素的喜悦本身就表明天主来临，而不是一个人的受孕。诗人祈求延长生命，也证明了他所呼求者的神性；因为我们渴望从天主那里获得生命和保存，而不是从人。确实，\Person{厄立特里亚女先知}这样呼求天主：\Quote{主啊，你为什么仍强迫我预告未来，而不将我从此地接走，以等待你来临的福日？}而\Person{玛罗}补充了他之前所说的：\par

开始吧，甜美孩童！你母亲微笑相迎，\newline{}她十个月来承受着你的重负。\newline{}没有凡人父母为你的诞生微笑：\newline{}你不知婚宴之乐，不属地之盛宴。\par

他的父母怎能向他微笑？因为他的父亲是天主，他是没有可感特质的德能，不以任何确定的形状存在，而是包含其他存有，因此并不在人身体内。谁不知道圣神不参与婚姻结合？因为在那众善皆向往的善之性情中，能有何欲望？总之，智慧与欢乐有何相通？但这些论证留给那些认为他有人类起源的人，以及那些不愿在言语和行为上净化自己一切邪恶的人。你，虔敬，我呼求你助我言语，你正是纯洁之律，万福中最可渴望的，至圣希望的导师，不朽的确然许诺！你，虔敬，和你，仁慈，我钦崇。我们这些获得你助佑的人，因你的医治能力欠你永恒的感恩。但那些因自身对你有天生的仇恨而失去你援助的群众，同样远离了天主自己，不知他们生命和存在的原因，以及一切不虔敬者的原因，都与对万有之主的天主当有的敬拜相连：因为世界本身属于他，以及其中的一切。\par

% structure: p0078 source=h2 target=section reason=relative-heading-rank; base=h2

\section{第二十二章　\Emph{皇帝感恩地将他的胜利和一切福佑归于基督；并谴责暴君马克西米努斯的行为，其迫害的暴行反增添了宗教的光荣。}}

你，虔敬，我将自己昌盛的原因和我现在所拥有的一切归于你。我所有努力的成功结局为此真理作证：英勇事迹、战争胜利、战胜敌人的凯旋。这真理连伟大的城市自己也欣然颂扬。那深受爱戴之城的人民也同意这同一情感，尽管他们一度被毫无根据的希望所欺骗，选了一位不配自己的统治者，一位迅速因他的狂妄行为而受罚的统治者。但现在，当我与你，虔敬，交谈，并认真努力以神圣温和的话语向你陈辞时，我还是不要回忆这些事吧。然而我要说一件事，或许并非不合时宜或不体面。暴君们曾对你，虔敬，和你的圣教会发动一场狂暴、残酷、无情的战争；而在罗马本身，也不乏有人幸灾乐祸于如此严重损害公共福祉的灾难。不，战场已备好；当你挺身而出，自愿作为献祭，以对天主的信德为支持时，那时，不虔敬之人的残酷，像吞噬之火般不停燃烧，却为你造就了奇妙而永世难忘的光荣。旁观者自己都惊愕不已，因为他们看见拷打神圣殉道者身体的刽子手自己也疲惫不堪，对暴行感到厌恶：锁链松脱，刑具无力，火焰熄灭，而受难者甚至一刻也未动摇其坚毅。那么，你们这些最不虔敬之人，通过这些残暴行径赢得了什么？你们疯狂愤怒的原因是什么？你肯定会说，这些行为是为敬拜诸神而做的。这些是什么神？或者你对神性有什么恰当的概念？你认为诸神像你一样易怒吗？如果真如此，你最好惊讶他们奇怪的决意，而不是服从他们严厉的命令，当他们催促你无辜屠杀义人时。你或许会引证祖宗的习俗和一般人的意见作为这种行为的理由。我承认事实：因为这些习俗与行为本身非常相似，出自同一愚蠢之源。你也许认为某种特殊力量存在于由人类艺术塑造和制造的偶像中；因此你尊敬并小心谨慎，唯恐它们被玷污：那些强大而崇高的神，竟如此依赖人的照顾！\par

% structure: p0080 source=h2 target=section reason=relative-heading-rank; base=h2

\section{第二十三章　\Emph{论基督徒的品行。天主喜悦那些过德性生活的人；我们必须期待审判和未来的报应。}}

将我们的宗教与你的比较。难道在我们这里没有真正的和谐和不懈的爱德吗？如果我们责备过错，我们的目的不是毁灭，而是劝诫；我们的纠正是为了安全，而非残忍？我们不仅对天主有真诚的信德，而且在社会生活的各种关系中也有忠信吗？我们不怜悯不幸者吗？我们的生活不是一种纯朴，不屑于以欺诈和伪善的面具掩盖邪恶吗？我们不承认真天主及其独一的统治权吗？这是真正的虔敬；这是真诚而真正无玷的宗教；这是智慧的生活；拥有这种生活的人，可说是在一条通往永生的崇高道路上旅行。因为走上这条道路并保持灵魂纯净、不受身体玷污的人，并非完全死亡；更确切地说，他完成了天主指派给他的服务，而非死亡。再者，承认效忠天主的人不会轻易被傲慢或暴怒压倒，而是在压力下高贵地站立，他的坚毅考验，可说是获享天主恩宠的通行证。因为我们不能怀疑神性悦纳人类行为的卓越。因为如果说有权势者和卑微者都感谢那些对他们有益的人，并以服务回报服务，而那位至高无上、统御万有者，甚至就是善本身，却在这事上疏忽，那将是荒谬的。相反，他贯穿我们生命的全程，在我们每次善行中靠近我们，接受并立即回报我们的德行和服从；尽管他将完全的酬报延迟到那未来的时期，那时我们一生的行为将受他审查，那些账目清白的人将获得永生的赏报，而恶人将遭受他们罪行应得的惩罚。\par

% structure: p0082 source=h2 target=section reason=relative-heading-rank; base=h2

\section{第二十四章　\Emph{论德西乌斯、瓦勒里安和奥勒良，他们因迫害教会而遭受悲惨结局。}}

我如今向你，\Person{德西乌斯}呼吁，你曾以侮辱践踏义人的劳苦；向你，教会的憎恨者，圣洁生活者的惩罚者：你现在死后情况如何？你现在的处境是多么艰难和悲惨！不仅如此，你死前的间隔已足够证明你悲惨的命运，当你在\Place{西徐亚平原}上连同全军覆没，使\Place{罗马}夸耀的武力受哥特人鄙视。还有你，\Person{瓦勒里安}，你对天主的仆人也表现出同样的残忍精神，你成了一个公义审判的榜样。你在敌人手中被俘，仍穿着紫袍和皇帝服饰时就被锁链牵着，最后你的皮肤被剥下，并由\Place{波斯}王\Person{沙普尔}下令保存，你留下了你灾难的永久纪念。还有你，\Person{奥勒良}，一切不义的凶猛行凶者，你的堕落多么显著！当你在\Place{色雷斯}疯狂的行程中，你在公路上被杀，不敬神的血灌满了道路的犁沟！\par

% structure: p0084 source=h2 target=section reason=relative-heading-rank; base=h2

\section{\Emph{第二十五章　\Person{戴克里先}可耻地退位，并因迫害教会而畏惧闪电的恐怖}}

\Person{戴克里先}作为迫害者展现了无情的残酷之后，却显示出对自己罪过的意识，并由于精神错乱的痛苦，忍受着卑微而独居的囚禁。那么，他积极敌对我们的天主获得了什么？我相信，就是这样：他在余生中持续畏惧闪电的打击。\Place{尼科美底亚}证明了这一事实；目击者，其中包括我自己，都这样宣告。宫殿和皇帝的私人卧室被摧毁，被闪电烧毁，被天火吞噬。有悟性的人确实预见了这种行为的后果；因为他们不能保持沉默，也不能掩饰对如此不配之事的悲痛，而是大胆公开地表达他们的感受，彼此说：\Quote{这是何等的疯狂！何等的傲慢滥用权力，竟敢与天主作战；故意侮辱最神圣、最正义的宗教；毫无理由地计划毁灭如此众多的义人？对他的臣民来说是节制的罕见榜样！对士兵来说是关怀和保护同胞公民的合格导师！从未见过败退之军后背的人，却将刀剑刺入同胞的胸膛！}流出的血如此之多，如果是在与野蛮敌人的战斗中流出，足以买到永久的和平。最后，确实，天主的眷顾对这些不圣洁的行为进行了报复，但国家也遭受了严重损失。因为我刚才所提皇帝的全部军队，受了一个在\Place{罗马}以暴力篡夺最高权力的无用之人的控制（当时天主的眷顾恢复了那座伟大城市的自由），在接连几次战斗中都被摧毁。当我们想起那些受压迫并热切渴望天生自由的人呼求天主帮助时的呼喊，以及当他们重新获得自由，恢复了自由平等的交往，在消除了他们曾呻吟的罪恶之后，对天主的赞美和感谢时：这些事情岂不在各方面都提供了令人信服的证据，证明天主的眷顾和他对人类福祉的慈爱关心吗？\par

% structure: p0086 source=h2 target=section reason=relative-heading-rank; base=h2

\section{\Emph{第二十六章　皇帝将其个人虔诚归于天主；并表明我们必须从天主那里寻求成功，并将成功归于他；但视为错误是我们自己疏忽的结果}}

当人们称赞我的服务时，这些服务源于上天的启示，他们岂不明显确立了天主是我所完成事业的原因这一真理吗？的确如此：因为属于天主的是做最好的事，而属于人的是履行天主的命令。我相信，最好最高贵的行动方式是，在尝试之前，我们尽可能为稳妥的结果做准备。而且所有人都知道，这些双手所从事的圣洁服务源于对天主的纯正真诚的信德；为公益所做的一切，都是通过积极努力，结合恳求和祈祷完成的；其结果是个体和公共福祉都达到了每个人敢于为自己和最亲爱的人所希望的程度。他们见证了战争，并目睹了天主的眷顾将胜利赐予这个人民的争斗；他们看见了他如何眷顾并回应我们的祈祷。因为义人的祈祷是不可战胜的；凡向天主献上圣洁祈求的人，没有不达到其目标的；除非信德动摇，否则不可能被拒绝；因为天主永远眷顾，永远乐意认可人类的德行。因此，人有时会犯错是自然的，但天主不是人错误的原因。因此，所有虔诚的人首先应为我们的个人安全，然后为公共事务的幸福状态，向万有的救主感恩；同时以圣洁的祈祷和不断的恳求恳求基督的恩宠，愿他继续赐予我们当前的祝福。因为他是义人不可战胜的盟友和保护者；他是万物的最高审判者，不朽之君，永生之赐予者。\par

\SourceRecord{Translated by Ernest Cushing Richardson.；Nicene and Post-Nicene Fathers, Second Series；Vol. 1.；Edited by Philip Schaff and Henry Wace.；Buffalo, NY: Christian Literature Publishing Co.,；1890.；<http://www.newadvent.org/fathers/2503.htm>.}
